\subsection{A relative two-simplex repair and its exact period obstruction}

We keep the notation $D(n)$ of AT13 and write
$K(1)=\operatorname{DK}(R[1])$ for the bottom submodule
of $D(2)$. The quotient is $K(2)=\operatorname{DK}(R[2])$.
The inclusions split in each simplex with free cokernel.

We first specify the elementary period map used below.
Fix a weight at which a target $C$ has its coefficient
comparison through the needed positive degrees.
A relative edge form means an element of $C_1$ whose
two faces are zero. For a closed relative one-form $\omega$,
choose $h\in C_1^0$ with $dh=\omega$. Its two faces are
coefficients, and define
\[
             \tau_C(\omega)=d_0h-d_1h\in S_R^w.
 \tag{AT.61}
\]
This is independent of $h$, since two primitives
differ by a coefficient in $C_1^0$, whose faces are equal.
Its kernel consists precisely of differentials of
relative zero-forms: when the two coefficient faces
are equal, subtract their common coefficient from $h$.
If the form-zero horizontal row is acyclic, every
coefficient is such a period. Indeed choose
$h\in N_1C^{w,0}$ with $d_0h=s$; then $dh$ is a
closed relative one-form of period $s$.
All these statements hold after rationalization as well,
and the period is natural under the specified
coefficient-preserving maps.

We also record two exact filling facts. Horizontal
acyclicity in a row makes $N_2C\to Z_1(NC)$ surjective
in that row. It also makes the endpoint map
$C_1\to C_0\oplus C_0$ surjective: fill the difference
of the two proposed endpoints by a normalized edge,
and add the degeneracy of the second endpoint.
Therefore a closed relative $j$-form, $j\geq2$, has a
relative primitive whenever the coefficient comparisons
through degree $j$ and the requisite horizontal row
bounds hold. First choose an absolute primitive in
$C_1^{j-1}$. Its two faces are closed positive forms
in $C_0$, hence differentials. Fill their primitives
by an edge and subtract its differential.
The resulting primitive has both faces zero.
These are constructions of actual cochain elements.

\begin{theorem}[AT15: Proved; a relative repair preserving simplices zero and one]
\label{AT15}
Let $X\to A,B$ be actual maps of strict simplicial
PD-DGAs over $R=\mathbb Z$ or $\mathbb Z_{(p)}$,
with torsionfree terms. Suppose $A,B$ are admissible,
$X$ has eventual horizontal acyclicity,
$H^0(X_a)=S_R$ in all simplices and weights, and
$S_{\mathbb Q}\to X_{a,\mathbb Q}$ is a quasi-isomorphism.
Suppose further that $X_0=S_R$ and that
$S_R\to X_1$ is an integral quasi-isomorphism
in every weight.

Fix $q\geq2$ and a weight $w\geq q$ above
\[
 \lambda_A(q),\ \lambda_B(q),\
 \mu_A(q),\ \mu_B(q),\ \mu_X(q).
 \tag{AT.62}
\]
Given any simplicial linear cycle map
$z:D(2)\to Z^q(X^{w,*})$, there are an integer
$m\ne0$ and actual simplicial maps
\[
 a:K(1)\to X^{w,q-1},\qquad
 b:D(2)\to X^{w,q-1},\qquad
 da=z|_{K(1)},\quad db=mz,\quad b|_{K(1)}=ma.
 \tag{AT.63}
\]
Form the relative two-cell algebra by adjoining
$u(D(2))$ in form degree $q-1$ and $v(D(2))$
in degree $q-2$, both of weight $w$, and imposing
\[
 du=z,\qquad dv=mu-b,\qquad
 u|_{K(1)}=a,\qquad v|_{K(1)}=0.
 \tag{AT.64}
\]
It is an actual strict PD extension $X\hookrightarrow Y$,
with $Y_0=X_0$ and $Y_1=X_1$, unchanged smaller weights,
and torsionfree terms, which remain free when the terms
of $X$ are free. Rationally the inclusion is a
quasi-isomorphism in every simplex and weight.
For $q=2$ and for every $q\geq4$, both maps to $A,B$
extend to $Y$; also $H^0(Y_a)=S_R$.
For $q\geq3$, its $H^1$ is unchanged.

For $q=3$, each target has the following exact obstruction
to extending to this specified relative algebra.
Choose a simplicial primitive $U$ of the image of $z$
with $U|_{K(1)}=fa$, and then choose any simplicial
primitive $V$ of $mU-fb$. Both choices exist.
The residue
\[
          \epsilon_f=\tau_A(V|_{K(1)})\bmod mS_R^w
                       \in S_R^w/mS_R^w
 \tag{AT.65}
\]
is independent of those choices. An actual map
$Y\to A$ extending $f$ exists exactly when
$\epsilon_f=0$. The same assertion, with its own
explicitly defined residue, holds for $B$.
\end{theorem}

\begin{proof}
Write $z_1=z|_{K(1)}$. It is a closed relative edge
form. Since $q\geq2$ and $X_1$ resolves $S_R$, it has
an integral primitive $a$ in $X_1^{w,q-1}$.
The two faces of this primitive are zero because
$X_0$ has no positive form degrees.
It therefore defines the indicated map from $K(1)$.
Projectivity of $D(2)$ and the rational coefficient
resolution of $X$ give a rational simplicial primitive
of $z$. Clearing the denominator of its one normalized
generator as in AT14 gives an initial $b$ and $m$
with $db=mz$. Its bottom differs from $ma$ by a
closed relative form
\[
                       c=b|_{K(1)}-ma.
 \tag{AT.66}
\]

For $q\geq3$ this form has degree at least two.
It has a primitive $h_1\in X_1^{w,q-2}$, and that
primitive is automatically relative since its degree
is positive and $X_0=S_R$.
The row bound in \textup{(AT.62)} gives a normalized
$h_2\in N_2X^{w,q-2}$ with $d_0h_2=h_1$.
It specifies a map $h:D(2)\to X^{w,q-2}$.
Replacing $b$ by $b-dh$ makes its bottom exactly
$ma$ while leaving $db=mz$.

For $q=2$ an integrality argument is required before
making that correction. Use either actual target map,
say $f:X\to A$. Since $D(2)$ is projective, eventual
vertical exactness gives a simplicial primitive
$U:D(2)\to A^{w,1}$ of $fz$.
The rational difference $U-fb/m$ is a closed
$D(2)$-map in form degree one. It is a differential
of a $D(2)$-map in degree zero, again by projectivity
and the rational vertical comparison.
The bottom of the latter map is a relative edge
zero-form, because $D(2)_0=0$.
Consequently its differential has period zero.
Naturality of \textup{(AT.61)} now gives
\[
 \frac{\tau_X(c)}m
   =\tau_{X_{\mathbb Q}}(b|_{K(1)}/m-a)
   =\tau_A(U|_{K(1)}-fa)\in S_R^w.
 \tag{AT.67}
\]
Thus this rational quotient is an integral coefficient.
By the form-zero row bound for $X$, choose a closed
relative one-form $r$ with period $\tau_X(c)/m$,
using a normalized coefficient edge and its vertical
differential as above.
Replace $a$ by $a+r$. Equation $da=z_1$ is unchanged
and the new $c$ has period zero.
It therefore has an integral relative zero-form
primitive $h_1$ in $X_1$, by the kernel description
of \textup{(AT.61)}.
Fill $h_1$ horizontally by $h_2\in N_2X^{w,0}$
and replace $b$ by $b-dh$.
This proves all three identities in \textup{(AT.63)}
also for $q=2$.

The relative algebra is now fully defined.
Take the free algebra with the first two differentials
in \textup{(AT.64)}, and quotient by the PD ideal
of the last two relations.
The differential preserves that ideal: on its first
generator it is $z_1-da=0$, and on the second it is
$ma-b|_{K(1)}=0$.
All relations are simplicial and homogeneous.
The split inclusion $K(1)_a\subset D(2)_a$ shows,
by the free PD universal property, that in each simplex
this is a free extension of $X_a$ on the two copies
of the free quotient $K(2)_a$.
This proves injectivity, the claims about terms and
weights, and literal equality in simplices zero and one.

Over $\mathbb Q$ change variables by $u'=u-b/m$,
keeping $v$ and $X$ fixed. The inverse is $u=u'+b/m$.
The boundary relations become $u'|_{K(1)}=0$,
$v|_{K(1)}=0$, and the differential becomes
$du'=0$, $dv=mu'$.
Thus this is the free identity disk on the actual
quotient parameter $K(2)$, over $X_{\mathbb Q}$.
Its simplexwise rational contractions are those
written in AT14. This proves the rational
quasi-isomorphism assertion.

We next construct the target extensions.
First choose any simplicial primitive
$U:D(2)\to A^{w,q-1}$ of $fz$.
The bottom difference $U|_{K(1)}-fa$ is a closed
relative $(q-1)$-form.
For $q\geq3$ it has a relative primitive by the
filling facts preceding the theorem.
For $q=2$ its period is zero: the now established
identity $b|_{K(1)}=ma$ and the rational comparison
of $U$ with $fb/m$ used in \textup{(AT.67)} give
that zero period. It too has an integral relative
primitive.
In either case fill this primitive horizontally
by a normalized two-simplex and subtract its vertical
differential from $U$.
This gives an actual simplicial primitive with
\[
                              U|_{K(1)}=fa.
 \tag{AT.68}
\]
The required horizontal fillings use only the
target row bounds through degree $q$ in
\textup{(AT.62)}.

The map $mU-fb$ is a closed $D(2)$-map of degree
$q-1\geq1$ and has zero bottom.
Projectivity and target exactness give a simplicial
primitive $V$ in degree $q-2$.
Its bottom is closed and relative.
For $q=2$ it is a closed relative zero-form, and
hence is already zero: it is a coefficient whose
two vertex values vanish.
For $q\geq4$ its degree is at least two, so it
has a relative primitive.
Fill that primitive horizontally and subtract
its differential from $V$, making the bottom zero.
We now have both required boundary values
\textup{(AT.68)} and $V|_{K(1)}=0$.
These actual generator values descend through
the defining relations of $Y$ and give $Y\to A$.
The construction for $B$ is the same, using the
same source data $a,b,m$.
For $q=2$, the actual map to $A$, torsionfreeness,
the rational coefficient comparison, and the weight
cutoff prove $H^0(Y_a)=S_R$ by the coefficient
integrality argument of AT13.
For $q\geq3$, all added generators have positive
form degree, so $H^0$ is unchanged directly.

It remains to prove the exact criterion when $q=3$.
For fixed $U$, two choices of $V$ differ by a closed
$D(2)$-map of form degree one.
It is a differential of a $D(2)$-map in degree zero.
Its bottom is a differential of a relative zero-form,
and thus has period zero.
For two choices of $U$ satisfying \textup{(AT.68)},
their difference is a closed $D(2)$-map of degree
two with zero bottom.
Choose a $D(2)$-primitive $\Omega$ of that difference
in degree one.
Its bottom is a closed relative one-form.
Subtracting $m\Omega$ from the resulting difference
of $V$'s leaves a closed $D(2)$-map in degree one,
whose bottom again has zero period.
Therefore the period of $V$ changes by an element
of $mS_R^w$, proving independence of
\textup{(AT.65)}.

Every such change is realizable.
Given $s\in S_R^w$, choose a closed relative one-form
$r$ in $A_1$ with period $s$.
Horizontal acyclicity in form degree one gives
$\Omega\in N_2A^{w,1}$ with $d_0\Omega=r$.
Regard it as a $D(2)$-map.
Its differential has zero bottom.
Replacing
\[
                     U\ \text{by}\ U+d\Omega,
                \qquad V\ \text{by}\ V+m\Omega
 \tag{AT.69}
\]
preserves all primitive equations and changes
the bottom period of $V$ by $ms$.
If \textup{(AT.65)} is zero, choose $s$ to make
that period zero.
Then the bottom of $V$ has a relative zero-form
primitive; fill it by a normalized two-simplex
and subtract its differential, making the bottom zero.
This constructs the desired relative map.
Conversely a relative map supplies choices with
bottom $V=0$, so its residue is zero.
This proves necessity and sufficiency.

Finally, for $q\geq4$ the new variables cannot
change $H^1$.
For $q=3$ the only new degree-one variables are
the $v$'s.
In each simplex the coefficient of every independent
new $u$ in the differential of a degree-one element
is $m$ times its corresponding $v$ coefficient.
Torsionfreeness forces those coefficients to vanish
on a cycle, exactly as in AT14.
The form-zero boundaries were unchanged.
This proves the asserted preservation of $H^1$.
\end{proof}

The residue in \textup{(AT.65)} is an obstruction
for the fixed boundary data in \textup{(AT.63)}.
In particular, the properties $X_0=S_R$ and
$X_1\simeq S_R$ alone do not assert its vanishing.
The scaled even disk on a $K(2)$ parameter examined
in the accompanying period-obstruction calculation
gives a nonzero example.

\paragraph{Horizontal restoration after an actual relative extension.}
Apply the construction of C21 starting from any $Y$
above for which both target maps have actually been
constructed. Use the fixed row thresholds
\[
 W_{\mathrm{hor}}(r)=
       \max\{1,r+1,\mu_A(r),\mu_B(r)\},
 \qquad
 \mu_{\mathrm{new}}(d)=\max_{r\leq d}W_{\mathrm{hor}}(r).
 \tag{AT.70}
\]
Omit zero cycles in positive-form homological degree
zero: that normalized term is already zero since
$Y_0=S_R$.
All other actual parameter fillings exist by the
target row bounds, exactly as in C21.
The countable repetitions at each form degree give
an algebra with the displayed uniform horizontal
bound and both extending maps.
It still has simplex zero $S_R$.
Only the form-zero homological-degree-zero cells
change simplex one, and those are free even-to-odd
vertical disks. Their integral contractions prove
that simplex one still resolves $S_R$ in every weight.
In each simplex every cell is a free vertical
identity-disk extension.
Its rational contraction preserves the rational
coefficient resolution.
The even-to-odd disks contract integrally; the
positive-form odd-to-even disks have no nonconstant
nonunit pair beginning below form degree three.
Expanding their tensor products as the direct sum
of constants and their displayed pairs therefore
shows that $H^0$ and $H^1$ are preserved integrally.
Filtered colimits are exact, so these assertions
persist through all repetitions and stages.
Terms stay torsionfree, and free when they were
free initially.
This constructs a horizontal restoration after the
specified individual relative repair; it asserts
no complete iteration of all vertical repairs.

